%! Parametric Functions — Unit 8, Lesson 8.1 — Warm-Up (Answer Key)
\documentclass[10pt]{article}
\usepackage{precalculus-article}
\usepackage{precalculus-key}   % pulls in -boxes; do NOT also load -boxes
\newcommand{\TallMath}[1]{$\displaystyle #1\rule[-1.4em]{0pt}{3.2em}$}

\begin{document}

\pageheader{Unit 8, Lesson 8.1}{Warm-Up --- Answer Key}
\namedateperiod

\vspace{0.12in}

\begin{notesbox}{Spiral Review --- Key}

\textbf{1.}\quad Complete the table of values for $f(x) = x^2 - 2x$.

\vspace{6pt}
\begin{center}
\renewcommand{\arraystretch}{1.5}
\begin{tabular}{|c|c|c|c|c|}
\hline
$x$ & $-1$ & $0$ & $1$ & $2$ \\
\hline
$f(x)$ & \ans{$3$} & \ans{$0$} & \ans{$-1$} & \ans{$0$} \\
\hline
\end{tabular}
\end{center}

\vspace{0.16in}
\textbf{2.}\quad Let $x(t) = 3t - 1$. Complete the table.

\vspace{6pt}
\begin{center}
\renewcommand{\arraystretch}{1.5}
\begin{tabular}{|c|c|c|c|c|}
\hline
$t$ & $0$ & $1$ & $2$ & $3$ \\
\hline
$x(t) = 3t - 1$ & \ans{$-1$} & \ans{$2$} & \ans{$5$} & \ans{$8$} \\
\hline
\end{tabular}
\end{center}

\vspace{0.16in}
\textbf{3.}\quad Let $x(t) = t + 1$ and $y(t) = t^2 - 3$.
Complete the table of ordered pairs $(x(t),\, y(t))$.

\vspace{6pt}
\begin{center}
\renewcommand{\arraystretch}{1.6}
\begin{tabular}{|c|c|c|c|}
\hline
$t$ & $x(t) = t + 1$ & $y(t) = t^2 - 3$ & $(x,\; y)$ \\
\hline
$-1$ & \ans{$0$} & \ans{$-2$} & \ans{$(0,\,-2)$} \\
\hline
$0$  & \ans{$1$} & \ans{$-3$} & \ans{$(1,\,-3)$} \\
\hline
$1$  & \ans{$2$} & \ans{$-2$} & \ans{$(2,\,-2)$} \\
\hline
\end{tabular}
\end{center}

\vspace{0.14in}
\textbf{4.}\quad Using your answers from Problem 3, which value of $t$
gives the \emph{lowest} $y$-value? How do you know?

\vspace{4pt}
\ans{$t = 0$} gives the lowest $y$-value of \ans{$-3$}. \ans{From the table, $y(0) = -3$ is smaller than $y(-1) = -2$ and $y(1) = -2$.}

\begin{teachernote}
  Students may notice that $y(t) = t^2 - 3$ is a parabola in $t$ with
  minimum at $t = 0$. Preview: this will be a key idea in Lesson 4.2
  when locating vertical extrema of particle motion.
\end{teachernote}

\end{notesbox}

\end{document}
